{{دیوان
| عنوان = مسولینی (کیا زمانے سے نرالا ہے مسولینی کا جُرم)
| ماخذ = ویکی ماخذ
| ماخذ_لنک = https://ur.wikisource.org/wiki/%D9%85%D8%B3%D9%88%D9%84%DB%8C%D9%86%DB%8C_%28%DA%A9%DB%8C%D8%A7_%D8%B2%D9%85%D8%A7%D9%86%DB%92_%D8%B3%DB%92_%D9%86%D8%B1%D8%A7%D9%84%D8%A7_%DB%81%DB%92_%D9%85%D8%B3%D9%88%D9%84%DB%8C%D9%86%DB%8C_%DA%A9%D8%A7_%D8%AC%D9%8F%D8%B1%D9%85%29
| مقابلہ = اقبال سائبر لائبریری (اقبال اکادمی پاکستان)
| مقابلہ_لنک = https://www.iqbalcyberlibrary.net
}}
=== (اپنے مشرقی اور مغربی حریفوں سے) ===

<poem>
کیا زمانے سے نرالا ہے مسولینی کا جُرم!
بے محل بِگڑا ہے معصومانِ یورپ کا مزاج

مَیں پھٹکتا ہُوں تو چھَلنی کو بُرا لگتا ہے کیوں
ہیں سبھی تہذیب کے اوزار! تُو چھَلنی، مَیں چھاج

میرے سودائے مُلوکیّت کو ٹھُکراتے ہو تم
تم نے کیا توڑے نہیں کمزور قوموں کے زُجاج؟

یہ عجائب شعبدے کس کی مُلوکیّت کے ہیں
راجدھانی ہے، مگر باقی نہ راجا ہے نہ راج

آلِ سِیزر چوبِ نَے کی آبیاری میں رہے
اور تم دُنیا کے بنجر بھی نہ چھوڑو بے خراج!

تم نے لُوٹے بے نوا صحرا نشینوں کے خیام
تم نے لُوٹی کِشتِ دہقاں، تم نے لُوٹے تخت و تاج

پردۂ تہذیب میں غارت گری، آدم کُشی
کل رَوا رکھّی تھی تم نے، مَیں رَوا رکھتا ہوں آج!
</poem>

٭ ۲۲ اگست ۱۹۳۵ء کو بھوپال (شیش محل) میں لکھے گئے
